\begin{houseviewbox}[核心结论]
AIDC 产业链不是只有 18 只股票。本报告现在分三层：第一层是 8 大板块、80 个子环节的全景产业链池；第二层是 58 只 A 股核心候选；第三层才是可发布目标价/公允价值组合。经过 2026 年 7 月 1 日扩展刷新后，可发布组合从原 18 只扩展到 56 只，其中新增 38 个扩展模型：13 个有明示券商目标价，24 个采用 AStock 自建公允价值，1 个采用 PS/SOTP 里程碑目标。仅 3 只因盈利或模型分母不足降级观察。当前 2026 年 6 月 30 日 11:30 原模型快照和 2026 年 7 月 1 日扩展快照下，\textbf{工业富联、沪电股份、润泽科技}的证据链相对更完整；光模块、光器件和高端 PCB 仍是最直接的利润池，但价格已显著反映 800G/1.6T 与 AI PCB 的景气预期。报告动作以“回调验证、市场支撑观察、高估值风险、观察降级”为主，不把主题热度当作买入理由。
\end{houseviewbox}

\begin{riskbox}[发布状态]
本版本估值已经接入可审计目标价/公允价值证据：56 个标的均有当前价、股本/市值、2026E 收入、净利、EPS、估值方法、Bear/Base/Bull、目标价或目标区间、隐含空间和来源路径。其中原 18 个标的使用 2026 年 6 月 30 日模型；新增 38 个来自 2026 年 7 月 1 日扩展核心候选模型。扩展模型中，13 个使用明示 broker/Street 目标价校准，24 个使用 AStock 自建公允价值且 Street 权重为 0，1 个使用 PS/SOTP 里程碑目标。剩余 3 个核心候选给出公司级处置并留在观察名单。
\end{riskbox}

\begin{exhibitbox}[投资委员会排序表]
\begin{longtable}{L{1.7cm}L{2.2cm}R{1.5cm}R{1.5cm}R{1.6cm}R{1.8cm}R{1.6cm}L{2.0cm}}
\toprule
\textbf{代码} & \textbf{公司} & \textbf{现价} & \textbf{2026E EPS} & \textbf{基础锚} & \textbf{综合目标} & \textbf{空间} & \textbf{动作}\\
\midrule
\endhead
300308 & 中际旭创 & 1289.74 & 22.52 & 1013.5 & 1054.7 & -18.2\% & 市场支撑观察\\
300502 & 新易盛 & 608.94 & 12.17 & 487.0 & 505.4 & -17.0\% & 市场支撑观察\\
601138 & 工业富联 & 72.85 & 2.21 & 75.1 & 65.1 & -10.6\% & 市场支撑观察\\
002463 & 沪电股份 & 154.74 & 2.58 & 90.4 & 101.3 & -34.5\% & 高估值风险\\
300442 & 润泽科技 & 78.67 & 1.57 & 47.0 & 50.6 & -35.7\% & 高估值风险\\
300394 & 天孚通信 & 309.68 & 2.75 & 151.4 & 182.8 & -41.0\% & 高估值风险\\
300476 & 胜宏科技 & 351.66 & 5.92 & 189.4 & 238.7 & -32.1\% & 高估值风险\\
600183 & 生益科技 & 177.84 & 2.00 & 60.0 & 85.8 & -51.8\% & 高估值风险\\
000938 & 紫光股份 & 28.86 & 1.20 & 29.9 & 28.5 & -1.2\% & 市场支撑观察\\
002837 & 英维克 & 79.62 & 0.64 & 22.3 & 48.1 & -39.6\% & 高估值风险\\
\bottomrule
\end{longtable}
\sourcenote{Sina 2026-06-30 11:30 行情快照；akshare 财务摘要；AStock 估值模型。空间=综合目标价/现价-1。}
\end{exhibitbox}

排序权重为：直接收入证据 35\%，财务交付 25\%，估值空间 20\%，市场情绪与流动性 10\%，风险/证据缺口 10\%。动作标签定义为：\textbf{核心关注}表示当前目标空间和证据质量均支持优先跟踪；\textbf{回调验证}表示公司质量较好但价格需要更好安全边际；\textbf{市场支撑观察}表示现价主要由情绪和流动性支撑，需等待财务或订单确认；\textbf{高估值风险}表示当前价格已显著透支模型可验证增长。

\begin{exhibitbox}[下一季度验证桥与上下行触发]
\begin{tabularx}{\textwidth}{L{3.0cm}X X}
\toprule
\textbf{环节} & \textbf{上行触发} & \textbf{下行触发}\\
\midrule
服务器/交换机 & AI 服务器、高速交换机收入继续高增，毛利率和现金流不恶化 & 增收不增利、存货/应收扩张、客户订单延后\\
光模块/光器件 & 800G/1.6T 出货、ASP 和毛利率同步验证 & ASP 下行、客户砍单、CPO/LPO 路线导致价值转移\\
AI PCB/CCL & 高多层、HDI、高速 CCL 收入占比继续提升 & 扩产、良率或原材料压力侵蚀毛利\\
供配电/液冷 & UPS/HVDC、CDU、冷板从认证进入批量交付 & 只有样机或认证，没有收入、验收和毛利确认\\
AIDC 运营 & 新增 MW、上架率、客户租约和经营现金流改善 & 高折旧、低利用率、电价或融资成本上行\\
\bottomrule
\end{tabularx}
\sourcenote{analysis/chain\_earnings\_bridge.md；analysis/valuation\_model.md；data/growth\_driver\_model.json。}
\end{exhibitbox}

第一层结论是链条真且很宽。完整 AIDC 链条覆盖算力芯片/存储、服务器零部件、网络/光通信、PCB/材料、供配电、液冷、数据中心建设运营和下游需求八大块。56 只目标价/公允价值标的只是当前公开财务数据和估值可复算性更完整的子集，不代表全产业链边界。

第二层结论是分化更真。直接兑现顺序大致为：服务器/交换机整机与代工、光互联、AI PCB/CCL、供配电和液冷、AIDC 运营商。越靠近最终算力服务，收入确认越依赖交付、上架率、电价和客户利用率；越靠近上游部件，估值越依赖订单能见度、ASP 与良率。

第三层结论是估值约束必须前置。多数核心标的午盘成交额和涨幅显示情绪处于强势状态；如果后续季度不能同时兑现收入增速、毛利率和现金流，当前价格隐含的增长久期会迅速收缩。
